\chapter{Arithmetic Co--Contra Comparison}
\label{v4-arith:arith-positselski}

An arithmetic co--contra comparison must specify its coalgebra,
continuous dual, completion, and derived categories over a common
coefficient ring. The classical field-coefficient correspondence
provides a model for this question. The module-theoretic
\(\mu\)-invariant does not characterize conilpotence, as shown in
Proposition~\ref{v4-arith:coefficients:prop:mu}.

For the proposed arithmetic Heisenberg application, a categorical
comparison must also preserve the relevant pairing and central
classes. Theorem~\ref{v4-arith:thm:P1-resolution} isolates the further
Mellin evaluation and defect-detection hypotheses that imply
\(\kappa+\kappa^!=0\Longleftrightarrow\xi(s)=\xi(1-s)\).
The scalar functional equation and the oscillator energy calculation
do not construct those arithmetic maps.

\section{Statement and Motivation}
\label{v4-arith:pos:statement}

\subsection{Review: Classical Positselski}
\label{v4-arith:pos:classical}

Let \(k\) be a field and \(C\) a dg coalgebra over \(k\).
A left \(C\)-comodule has a coaction \(M\to C\otimes_k M\).
A left \(C\)-contramodule has a contraaction
\(\pi:\operatorname{Hom}_k(C,P)\to P\).
The latter satisfies the associativity identity induced by
\(\Delta:C\to C\otimes_k C\) and the unit identity induced by
\(\epsilon:C\to k\). All maps commute with differentials,
and tensor--Hom adjunctions use the Koszul signs.

The \emph{coacyclic} dg comodules form the smallest triangulated
subcategory containing totalizations of exact triples and closed
under arbitrary direct sums. The \emph{contraacyclic} dg
contramodules are defined using products instead.
Their Verdier quotients are the coderived and contraderived
categories. These quotients need not agree with the ordinary
derived categories, even though the coefficient ring is a field.

For a right comodule \(N\) and a left contramodule \(P\), define
\[
 N\odot_C P=
 \operatorname{coker}\bigl(
 N\otimes_k\operatorname{Hom}_k(C,P)
 \rightrightarrows N\otimes_k P\bigr).
\]
The arrows use the contraaction and the right coaction followed
by evaluation, respectively. On homogeneous elements the latter
evaluation is \(c\otimes f\mapsto(-1)^{|c||f|}f(c)\).
This quotient is the \emph{contratensor product}.
Write \(\operatorname{Hom}_C\) for colinear homogeneous maps,
with their Hom differential.

\begin{theorem}[Positselski, \emph{Two kinds of derived categories}, \S5.2]
\label{v4-arith:pos:thm:classical-positselski}
The derived functors of
\[
 \Psi_C(M)=\operatorname{Hom}_C(C,M),\qquad
 \Phi_C(P)=C\odot_C P
\]
give mutually inverse equivalences
\begin{equation}
\label{v4-arith:pos:eq:classical}
 \mathrm D^{\mathrm{co}}(C\text{-comod}_k)
 \simeq
 \mathrm D^{\mathrm{ctr}}(C\text{-contra}_k).
\end{equation}
The functor \(\Phi_C\) is left adjoint to \(\Psi_C\).
\end{theorem}
\begin{proof}
Positselski's resolution theorem, \S4.4, represents the two
quotients by dg objects that are injective or projective after
forgetting the differential. Graded injective comodules are
summands of cofree objects \(C\otimes_k V\). Graded projective
contramodules are summands of free objects
\(\operatorname{Hom}_k(C,V)\).
The two functors exchange these free and cofree objects, with
natural inverse maps. They therefore exchange their summands
and the corresponding dg objects. The resulting equivalence
of these homotopy categories proves the assertion.
\end{proof}

This correspondence is covariant in \(M\).
For \(C=k\), its right adjoint is \(M\mapsto M\), not the
contravariant vector-space dual. The coalgebra dual \(C^\vee\)
and the cobar algebra \(\Omega C\) are also different constructions.
A topological description using \(C^\vee\) must specify its
continuous duality and its contramodules.

\subsection{The \texorpdfstring{\(\Lambda\)}{Lambda}-Coefficient Target}
\label{v4-arith:pos:lambda-target}

Fix a prime \(p\), a procyclic group \(\Gamma\cong\mathbb Z_p\),
and a topological generator \(\gamma\). Set
\[
 \Lambda=\mathbb Z_p[[\Gamma]]\cong\mathbb Z_p[[T]],
 \qquad T=\gamma-1,\qquad \mathfrak m=(p,T).
\]
For the cyclotomic application, \(\Gamma\) is the Galois group
of a chosen cyclotomic \(\mathbb Z_p\)-extension.
The full cyclotomic Galois group may also have a finite factor.
The ring \(\Lambda\) is complete, Noetherian, local and regular
of dimension two, with residue field \(\mathbb F_p\).

A pseudo-isomorphism is a map with finite kernel and cokernel.
Every finitely generated \(\Lambda\)-module \(M\) admits one
of the form
\begin{equation}
\label{v4-arith:pos:eq:lambda-structure}
 M\longrightarrow\Lambda^r
 \oplus\bigoplus_i\Lambda/p^{a_i}\Lambda
 \oplus\bigoplus_j\Lambda/f_j^{b_j}\Lambda.
\end{equation}
Here \(r\ge0\), \(a_i,b_j\ge1\), and each \(f_j\in\mathbb Z_p[T]\)
is irreducible and distinguished.
A monic polynomial is distinguished when its other coefficients
are divisible by \(p\).
The Iwasawa invariants are
\(\mu(M)=\sum_i a_i\) and
\(\lambda(M)=\sum_j b_j\deg f_j\).
Proposition~\ref{v4-arith:pos:prop:lambda-structure}
recalls the structure theorem.

\begin{definition}[arithmetic dg \(\Lambda\)-coalgebra]
\label{v4-arith:pos:def:arith-dg-coalgebra}
An \emph{arithmetic dg \(\Lambda\)-coalgebra} is a coaugmented
dg coalgebra \(C=\Lambda\mathbf1\oplus\overline C\) over
\(\Lambda\), with \(\overline C=\ker\epsilon\), satisfying:
\begin{enumerate}
\item \emph{Local finiteness over \(\Lambda\)}: the underlying
graded \(\Lambda\)-module \(C\) is a filtered union of finitely
generated \(\Lambda\)-submodules \(C_{\leq n}\) with
\(C=\varinjlim_{n}C_{\leq n}\) and each \(C_{\leq n}/C_{\leq n-1}\)
finitely generated.
\item \emph{Comultiplication}: a graded \(\Lambda\)-linear map
\(\Delta\colon C\to C\otimes_{\Lambda}C\) with counit
\(\epsilon\colon C\to\Lambda\), satisfying the standard
coassociativity / counit axioms, and compatible with the
differential \(d\).
\item \emph{Conilpotence}: for every \(c\in\overline C\), some
iterated reduced coproduct vanishes on \(c\). The reduced coproduct is
\[
 \overline\Delta(c)=\Delta(c)-\mathbf1\otimes c-c\otimes\mathbf1
 \in\overline C\otimes_\Lambda\overline C.
\]
\end{enumerate}
The notation \(\mu(C)=\mu(\overline C)\) requires
\(\overline C\) to be a finitely generated torsion
\(\Lambda\)-module. For an infinite graded or completed object,
one must instead specify the coefficient modules whose invariants
are being taken.
\end{definition}

\begin{definition}[\(\mu=0\) arithmetic \(\Lambda\)-coalgebra]
\label{v4-arith:pos:def:mu-zero}
Suppose \(M=\overline C\) is a finitely generated torsion
\(\Lambda\)-module. Then \(C\) is \emph{\(\mu=0\)} when
\[
 \mu(M)=\operatorname{length}_{\Lambda_{(p)}}M_{(p)}=0.
\]
Equivalently, its structure theorem has no
\(\Lambda/p^a\Lambda\) summand. This is a statement up to
pseudo-isomorphism, not the absence of all \(p\)-torsion.
For example, \(M=\Lambda/(p,T)\) is nonzero and killed by \(p\),
but \(M_{(p)}=0\) because \(T\) becomes invertible.
\end{definition}

\begin{remark}[the role of \(\mu=0\)]
\label{v4-arith:pos:rem:why-mu-zero}
The invariant measures torsion at the height-one prime \((p)\).
It does not control all torsion or derived base change.
The module \(\Lambda/(p,T)\) has \(\mu=0\) and nonzero
\(p\)-torsion. The module \(\Lambda/(T)\) also has \(\mu=0\),
but its first Tor with itself is \(\Lambda/(T)\).
Proposition~\ref{v4-arith:pos:prop:resolutions-cohere} computes
these effects. Completion and reduction must therefore be
checked on the actual comparison maps.
\end{remark}

A precise lifting criterion uses completion inside the recipient
category. In a stable \(\Lambda\)-linear \(\infty\)-category with
sequential limits, put
\[
 X/p^n=\operatorname{cofib}(p^n:X\to X).
\]
Call \(X\) \emph{derived \(p\)-complete} when the natural map
\(X\to\varprojlim_n X/p^n\) is an equivalence.
For \(X\in\mathrm D(\Lambda)\), the quotient is
\(X\otimes_\Lambda^{\mathbb L}\Lambda/p^n\).
In a coderived recipient, its intrinsic cofiber and limit are
required. An underlying ordinary derived complex may forget
information.

\begin{theorem}[lifting an arithmetic comparison]
\label{v4-arith:thm:arithmetic-positselski}
Let \(\mathcal C\) and \(\mathcal D\) be specified stable
\(\Lambda\)-linear \(\infty\)-categories with sequential limits.
Suppose an exact \(\Lambda\)-linear adjunction
\(\Phi:\mathcal D\rightleftarrows\mathcal C:\Psi\) has been
constructed. Write its unit and counit as
\[
 \eta_P:P\longrightarrow\Psi\Phi P,
 \qquad
 \varepsilon_M:\Phi\Psi M\longrightarrow M.
\]
Assume their cofibers are derived \(p\)-complete in their
respective categories, and that \(\eta_P/p\) and
\(\varepsilon_M/p\) are equivalences for every \(P,M\).
Then
\begin{equation}
\label{v4-arith:pos:eq:arith-positselski}
 \Psi:\mathcal C\xrightarrow{\ \sim\ }\mathcal D,
 \qquad \Phi\simeq\Psi^{-1}.
\end{equation}
\end{theorem}

The proof is given in Section~\ref{v4-arith:pos:assembly}.
For the arithmetic problem, one must construct \(\mathcal C\)
and \(\mathcal D\) from the intended integral comodules and
contramodules, with their exact structures and topology.
One must then construct the adjunction and verify these
hypotheses. Module-theoretic \(\mu=0\) does not supply them.
Compatibility with the algebraic bar--cobar adjunction is a
further comparison of functors.

\begin{remark}[data for the arithmetic comparison]
\label{v4-arith:pos:rem:open-sub-items}
The categorical construction and its zeta interpretation require
different data:
\begin{enumerate}
\item \emph{Continuous duals.} Specify the coalgebra, its topology,
its dual recipient, and the exact comodule and contramodule
categories. Identify any proposed topological-algebra description
with these categories.
\item \emph{Derived completion.} Prove completeness of the unit
and counit cofibers in those recipients. Ordinary module
completion need not detect coderived equivalences.
\item \emph{Reduction of the adjunction.} Construct the comparison
with an actual adjunction over \(\Lambda/p\), preserving its
unit and counit. Prove that the reduced maps are equivalences.
This requires detection of equivalences, not faithfulness of
reduction on morphisms.
\item \emph{Character and pairing.} Identify the selected
one-mode character with \(\zeta\) and construct its arithmetic
pairing. The full oscillator character is
\(\prod_{k\ge1}\zeta(k\beta)\). Neither scalar character
constructs the central-defect or Mellin maps.
\item \emph{Bost--Connes comparison.} Construct the cobar algebra
in the chosen coefficient category and compare it with the
specified Bost--Connes algebra, action and pairing.
Question~\ref{v4-arith:cor:koszul-dual-BC} retains this problem.
The continuous dual is not the cobar algebra by definition.
\end{enumerate}
The first three requirements concern the lifting criterion.
The last two concern its arithmetic interpretation.
They are not an exhaustion theorem for all possible arithmetic
realizations.
\end{remark}

Section~\ref{v4-arith:pos:main-theorem} presents the proposed
co--contra construction. Its six coefficient and completion
comparisons are in Section~\ref{v4-arith:pos:obstruction-analysis-table}.
Section~\ref{v4-arith:pos:iwasawa-preliminaries} fixes the required
\(\Lambda\)-module notions.

The arithmetic Heisenberg application and its level comparison
are in Section~\ref{v4-arith:pos:heisenberg-application}.
Section~\ref{v4-arith:pos:residual-domain} retains the non-abelian
Iwahori and Whittaker comparison conditions.

\section{Obstruction Analysis Table}
\label{v4-arith:pos:obstruction-analysis-table}

The following comparisons specify the coefficient, completion and
pairing conditions in the arithmetic statement.

\subsection{Obstruction 1: Base-Change Lift}
\label{v4-arith:pos:att-1-base-change}

The field correspondence is not a theorem over
\(\mathbb F_p[[T]]\), which is a discrete valuation ring.
There is no unital map from a field to \(\mathbb Z_p[[T]]\):
positive characteristic is incompatible with characteristic zero,
and a characteristic-zero field would make \(p\) invertible.
Thus extension of scalars from a field cannot construct the
integral comparison directly.

Reduction \(\Lambda\to\Lambda/p\) is nevertheless useful.
On derived \(p\)-complete objects, it detects equivalences,
by Proposition~\ref{v4-arith:pos:prop:lifting}.
This proves a lifting criterion once the functors and their
reduced unit and counit are available.
It does not construct these functors, and reduction is not
faithful on morphisms.

\subsection{Obstruction 2: Exact Structures and Derived Recipients}
\label{v4-arith:pos:att-2-co-contra}

Over \(\Lambda\), an arithmetic comparison must select the
admissible exact sequences of comodules and contramodules.
The coacyclic and contraacyclic subcategories are then formed
using these sequences and the appropriate sums or products.
An injective resolution over the coefficient ring need not be
an injective resolution in the comodule category.
The corresponding distinction holds for projectives.

Selmer complexes furnish arithmetic complexes with specified
local conditions. To use them here requires functors into the
chosen comodule or contramodule recipient, preserving those
conditions. Their existence does not follow from the notation
\(\mathrm D(\Lambda)\), or from a vanishing \(\mu\)-invariant.
Definitions~\ref{v4-arith:pos:def:co-derived} and
\ref{v4-arith:pos:def:contra-derived} state the exact-category
construction used below.

\subsection{Obstruction 3: \texorpdfstring{\(\mu\)}{mu}-Invariant Vanishing for the Heisenberg Bar}
\label{v4-arith:pos:att-3-mu-heisenberg}

For a defined arithmetic bar object, the first question is which
finitely generated torsion \(\Lambda\)-module carries the proposed
\(\mu\)-invariant. A filtered or completed complex has no single
such invariant until the relevant cohomology or graded coefficient
modules are specified.

There is an elementary vanishing calculation on a fixed coefficient
model. Let \(\Lambda=\mathbb Z_p[[T]]\), and let \(M\) be a finite
free \(\mathbb Z_p\)-module of rank \(d\), with trivial
\(\Gamma\)-action. Then
\[
 M\cong(\Lambda/(T))^d,
 \qquad M_{(p)}=0,
 \qquad \mu(M)=0.
\]
Indeed, \(T\notin(p)\), so localization at \((p)\) makes the
annihilator \(T\) invertible. This proves vanishing for this module
without identifying it with a cyclotomic Iwasawa module.

Given an actual augmented coefficient algebra with augmentation ideal
\(I\), one can form finite bar-length pieces
\[
 B_{\leq n}=\bigoplus_{j=0}^n(I[1])^{\otimes_\Lambda j},
\]
with cohomological shift \(I[1]^q=I^{q+1}\).
For positive bar lengths, applying that calculation requires the
selected coefficients to be finite free over \(\mathbb Z_p\) with
trivial action. The length-zero coefficient is \(\Lambda\) itself. A residue
complex can be tensored into this construction only after its
\(\Lambda\)-module structure and total differential have been supplied.
Finite tame symbols, by themselves, provide neither structure.

Ferrero--Washington gives a vanishing theorem for specified
cyclotomic Iwasawa modules. Its use requires identifying the arithmetic coefficient module
with one in that theorem's domain. The identification must preserve
the action and topology. Neither the name Heisenberg nor commutativity of
\(\Gamma\) supplies such an identification.
Proposition~\ref{v4-arith:coefficients:prop:mu} further separates
module-theoretic vanishing from coalgebra conilpotence.

\subsection{Obstruction 4: Weight Compatibility Under Bar--Cobar}
\label{v4-arith:pos:att-4-weights}

\paragraph{Weight comparison.} The proposed arithmetic pairing
would relate bar coefficients of energy \(\log p\) to independently
defined dual coefficients of energy \(-\log n\).
To interpret energy cancellation as
\(\kappa^{\mathrm{Ar}}+(\kappa^!)^{\mathrm{Ar}}=0\), one must
construct the central-defect identity and evaluation of
Theorem~\ref{v4-arith:thm:P1-resolution}. Compatibility with a
Bost--Connes Hecke operator \(\mu_n\) is an additional action map.

A cyclotomic character supplies a Galois action only after that
action is specified on the coefficient object.
Its compatibility with a co--contra functor is additional data.
Oscillator energies \(k\log p\) belong to a different grading.
The ungraded correspondence does not identify the two.

\paragraph{Weight calculation.}
Convention~\ref{v4-arith:conv:two-grading}(ii) assigns energy
\(k\log p\) to \(x_{p,k}=\alpha^{(p)}_{-k}\). Its linear dual can be
given the negative energy. Evaluation then has total energy zero,
as constructed in Theorem~\ref{v4-arith:pos:thm:weight-comp}.
The numerical identity \(\log(p^k)=k\log p\) is also compatible with
assigning \(-\log(p^k)\) to a proposed dual symbol. It does not identify
that symbol with a Bost--Connes isometry, construct a twisting cochain,
or identify oscillator energy with a central level.

The full one-species-per-prime oscillator trace is
\(\prod_{k\ge1}\zeta(k\beta)\) for \(\beta>1\). The selected mode-one
occupation space has trace \(\zeta(\beta)\), but is not invariant under
the divisor or index-shift operations for index greater than one
(Section~\ref{v4-arith:ado:hilbert}). A pairing or arithmetic comparison must
specify which of these carriers it uses.

\subsection{Obstruction 5: Archimedean Completion}
\label{v4-arith:pos:att-5-archimedean}

The completed zeta function contains the real factor
\(\pi^{-s/2}\Gamma(s/2)\). Fourier self-duality of the Gaussian
and Poisson summation provide its analytic input.
Theorem~\ref{v4-arith:thm:arithmetic-modular-invariance}
proves the scalar reflection with its polar terms.

A categorical interpretation needs an archimedean coalgebra or
algebra, its modules and its pairing. It also needs a
comparison to the Fourier model.
The local Fock space does not specify these data by itself.
Section~\ref{v4-arith:pos:archimedean} gives the Gaussian
calculation and a conditional assembly of actual local
equivalences over a common coefficient ring.

\subsection{Obstruction 6: Noncommutative Coefficients}
\label{v4-arith:pos:att-6-noncommutative}

For a compact \(p\)-adic group \(G\), the completed group algebra
\(\mathbb Z_p[[G]]\) is a possible noncommutative coefficient ring.
Its use for an Iwahori Hecke category requires an actual action
and a comparison with that category. These are not consequences
of the name of the group.

A coalgebra over a noncommutative base must specify left and
right module structures. The resulting coring has a coproduct
into its tensor square over the base.
Its conilpotence, continuous Hom, completion and exact structures
must be stated in that setting. The commutative structure theorem
below does not provide them. The noncommutative and Whittaker
questions in Section~\ref{v4-arith:pos:residual-domain} retain
these additional constructions.

\subsection{Comparison Table}
\label{v4-arith:pos:table-summary}

\begin{table}[h]
\centering
\small
\begin{tabular}{@{}p{.10\linewidth}|p{.34\linewidth}|p{.47\linewidth}@{}}
Input & Content & Conclusion \\
\hline
1 & Reduction and lifting & complete cofibers and actual maps \\
2 & Derived recipients & specify exact structures \\
3 & \(\mu=0\) for specified bar modules & requires the module comparison \\
4 & Weight pairing & oscillator evaluation constructed \\
5 & Archimedean comparison & Fourier model. Categorical map needed \\
6 & Noncommutative coefficients & coring and action required \\
\end{tabular}
\caption{Arithmetic Positselski comparisons. The abelian comparison retains
 the data of Remark~\ref{v4-arith:pos:rem:open-sub-items}
 and the pairing and invariant maps of
 Question~\ref{v4-arith:pos:thm:P1-full-unconditional}.
 Non-abelian Iwahori level requires further data
 (Section~\ref{v4-arith:pos:non-abelian}).}
\label{v4-arith:pos:tab:table}
\end{table}

\section{Iwasawa Preliminaries}
\label{v4-arith:pos:iwasawa-preliminaries}

The coefficient calculations use three facts about finitely
generated \(\Lambda\)-modules:
the structure theorem, the resolution bounds, and the actual
Tor groups governing reduction. These module calculations do not
identify the derived comodule and contramodule recipients.

\subsection{\texorpdfstring{\(\Lambda\)}{Lambda}-Modules and Structure Theorem}
\label{v4-arith:pos:lambda-modules}

\begin{proposition}[\(\Lambda\)-module structure theorem]
\label{v4-arith:pos:prop:lambda-structure}
Let \(\Lambda=\mathbb{Z}_{p}[[\Gamma]]\)
and let \(M\) be a finitely generated \(\Lambda\)-module. Then
there is a pseudo-isomorphism (morphism with finite kernel and
cokernel)
\begin{equation}
M\sim \Lambda^{r}\oplus\bigoplus_{i=1}^{s}\Lambda/p^{a_{i}}\Lambda\oplus\bigoplus_{j=1}^{t}\Lambda/f_{j}^{b_{j}}\Lambda,
\end{equation}
where \(r,s,t\in\mathbb{Z}_{\geq 0}\), \(a_{i}\in\mathbb{Z}_{\geq 1}\),
and \(f_{j}\in\mathbb{Z}_{p}[\gamma-1]\) are distinguished
irreducible polynomials (Weierstrass polynomials). The invariants
\(r\), \(\mu(M)=\sum a_{i}\), \(\lambda(M)=\sum b_{j}\deg f_{j}\) are
well-defined.
\end{proposition}

\begin{proof}
Apply Serre, \emph{Classes des corps cyclotomiques}, \S6,
Lemmas~5--6. Up to finite kernel and cokernel, a finitely
generated module splits into a reflexive module and height-one
primary factors. The reflexive module is free because \(\Lambda\)
is regular local of dimension two. Localization at a height-one
prime is a discrete valuation ring, so each primary factor has
its elementary-divisor decomposition. Weierstrass preparation
identifies the height-one primes as \((p)\) and the ideals
\((f)\) for distinguished irreducible \(f\).
Their localized multiplicities give the stated invariants.
\end{proof}

\begin{corollary}[\(\mu=0\) criterion]
\label{v4-arith:pos:cor:mu-zero-criterion}
For a finitely generated \(\Lambda\)-module \(M\), the invariant
\(\mu(M)\) is the length of the torsion submodule of \(M_{(p)}\).
It vanishes exactly when \(M_{(p)}\) is free over
\(\Lambda_{(p)}\). If \(M\) is torsion, this is equivalent
to \(M_{(p)}=0\).
\end{corollary}
\begin{proof}
Finite modules vanish after localization at \((p)\).
The structure theorem therefore becomes an isomorphism there.
The summand \(\Lambda/p^a\) contributes length \(a\).
A distinguished polynomial reduces modulo \(p\) to a nonzero
power of \(T\), so it does not lie in \((p)\) and becomes a
unit. Its summand vanishes. The remaining free summand has
rank \(r\), proving all assertions.
\end{proof}

\subsection{Injective and Projective Resolutions Over \texorpdfstring{\(\Lambda\)}{Lambda}}
\label{v4-arith:pos:resolutions}

\begin{proposition}[projective resolution bound]
\label{v4-arith:pos:prop:proj-resolution}
Every finitely generated \(\Lambda\)-module \(M\) has a finite
free resolution of length at most two.
\end{proposition}
\begin{proof}
The regular sequence \((p,T)\) gives the Koszul resolution
\[
 0\longrightarrow\Lambda
 \xrightarrow{\binom{-T}{p}}\Lambda^2
 \xrightarrow{(p\ \ T)}\Lambda
 \longrightarrow\mathbb F_p\longrightarrow0.
\]
It follows that \(\operatorname{Tor}_i^\Lambda(\mathbb F_p,M)=0\)
for \(i>2\). Construct a minimal free resolution of \(M\)
by successively choosing minimal generators of its kernels.
Each free module is finite because \(\Lambda\) is Noetherian.
Every differential reduces to zero modulo \(\mathfrak m\),
so its \(i\)-th free module modulo \(\mathfrak m\) is the
\(i\)-th Tor group. Nakayama's lemma makes those free modules
zero for \(i>2\).
\end{proof}

\begin{proposition}[injective resolution bound]
\label{v4-arith:pos:prop:inj-resolution}
Every \(\Lambda\)-module \(M\) has an injective resolution
\(0\to M\to I^0\to I^1\to I^2\to0\).
No finite-generation condition on the injective terms is asserted.
\end{proposition}
\begin{proof}
For every ideal \(J\), the preceding proposition gives
\(\operatorname{Ext}^3_\Lambda(\Lambda/J,M)=0\).
Start an injective resolution and let \(Q\) be the cokernel
of its first two injective terms, so that
\(0\to M\to I^0\to I^1\to Q\to0\) is exact.
Dimension shifting gives
\(\operatorname{Ext}^1_\Lambda(\Lambda/J,Q)=0\) for every \(J\).
Baer's criterion makes \(Q\) injective.
\end{proof}

These are resolutions of modules over \(\Lambda\).
They do not supply the required injectives or projectives in
an arithmetic comodule or contramodule category.
In particular, the residue field is not an injective
\(\Lambda\)-module: the same Koszul complex gives
\(\operatorname{Ext}^1_\Lambda(\mathbb F_p,\mathbb F_p)
\cong\mathbb F_p^2\).

\begin{proposition}[derived reduction and its torsion]
\label{v4-arith:pos:prop:resolutions-cohere}
Put \(B=\Lambda/p=\mathbb F_p[[T]]\).
For every \(\Lambda\)-module \(M\),
\[
 M\otimes_\Lambda^{\mathbb L}B
 \simeq[\,M\xrightarrow{p}M\,],
 \qquad
 \operatorname{Tor}_1^\Lambda(M,B)=M[p],
\]
where the displayed complex lies in degrees \(-1,0\), and
\(M[p]=\ker(p:M\to M)\). Higher Tor groups vanish.
Reduction of a projective resolution therefore resolves the
ordinary module \(M/pM\) exactly when \(M[p]=0\).
The corresponding statements for \(\Lambda/(T)\) use the
map \(T\) and its kernel \(M[T]\).
\end{proposition}
\begin{proof}
Tensor the free resolution
\(0\to\Lambda\xrightarrow{p}\Lambda\to B\to0\) with \(M\).
Its homology gives the assertion. Multiplication by \(T\) is
also injective on \(\Lambda\), so the same argument applies
to \(\Lambda/(T)\).
\end{proof}

For \(M=\Lambda/(p,T)\), one has \(\mu(M)=0\) and
\(M[p]=M\ne0\). Thus the positive-degree Tor term survives.
For \(N=\Lambda/(T)\), one has \(\mu(N)=0\) and
\[
 \operatorname{Tor}_1^\Lambda(N,N)=N\ne0.
\]
This second module is free over \(\mathbb Z_p\) but is not
flat over \(\Lambda\). Neither \(\mu=0\) nor
\(\mathbb Z_p\)-freeness can replace the required Tor calculation.

\subsection{Co-Derived and Contra-Derived \texorpdfstring{\(\Lambda\)}{Lambda}-Categories}
\label{v4-arith:pos:co-contra-lambda}

Let \(C\) be a specified dg \(\Lambda\)-coalgebra.
Fix exact structures on its dg comodules and dg contramodules,
including the admissible exact triples.
For example, one may use triples split as graded
\(\Lambda\)-modules. This gives a relative exact structure.
An abelian or continuous version requires its own hypotheses.

Assume the chosen comodule category has admissible arbitrary
sums and the contramodule category has admissible arbitrary
products. Write \(\mathrm{Hot}\) for the homotopy category
of these dg objects.

\begin{definition}[\(\Lambda\)-coderived recipient]
\label{v4-arith:pos:def:co-derived}
Let \(\mathrm{Ac}^{\mathrm{co}}_\Lambda(C)\) be the smallest
thick triangulated subcategory of
\(\mathrm{Hot}(C\text{-comod}_\Lambda)\) containing the
totalizations of admissible exact triples and closed under
arbitrary sums. Define
\[
 \mathrm D^{\mathrm{co}}_\Lambda(C\text{-comod})
 =\mathrm{Hot}(C\text{-comod}_\Lambda)
       /\mathrm{Ac}^{\mathrm{co}}_\Lambda(C).
\]
The chosen exact structure is part of this notation.
\end{definition}

\begin{definition}[\(\Lambda\)-contraderived recipient]
\label{v4-arith:pos:def:contra-derived}
Let \(\mathrm{Ac}^{\mathrm{ctr}}_\Lambda(C)\) be the smallest
thick triangulated subcategory of
\(\mathrm{Hot}(C\text{-contra}_\Lambda)\) containing the
totalizations of admissible exact triples and closed under
arbitrary products. Define
\[
 \mathrm D^{\mathrm{ctr}}_\Lambda(C\text{-contra})
 =\mathrm{Hot}(C\text{-contra}_\Lambda)
       /\mathrm{Ac}^{\mathrm{ctr}}_\Lambda(C).
\]
A description by continuous \(C^\vee\)-contramodules requires
an identification with this specified recipient.
\end{definition}

\begin{remark}[comparison with Selmer complexes]
\label{v4-arith:pos:rem:nekovar}
A Selmer complex is an object of an ordinary derived category
with specified local-condition maps.
Chapter~\ref{v4-arith:coefficients:chapter} constructs such
complexes and their coefficient comparisons.
A realization in the recipients above must preserve those maps
and specify which acyclic objects are annihilated.
The distinction between acyclic and coacyclic objects already
occurs for coalgebras over a field. Coefficient linear algebra
does not remove it.
\end{remark}

\begin{question}[enhancements and compact objects]
\label{v4-arith:pos:lem:representability}
For the intended arithmetic coalgebra and exact structures,
construct the stable enhancements of these recipients.
Determine the complete subcategories on which the co--contra
functors exist and the compact objects, if the relevant
coproducts exist.
\end{question}

Compactness cannot mean concentration in degree zero.
Compact objects in a stable category are closed under shifts.
For example, \(\Lambda[1]\) is compact in
\(\mathrm D(\Lambda)\) and is not a module concentrated in
degree zero. The module resolution bounds do not resolve
Question~\ref{v4-arith:pos:lem:representability}.

\subsection{Iwasawa \texorpdfstring{\(\mu\)}{mu}-Invariant and Ferrero--Washington}
\label{v4-arith:pos:mu-FW}

\begin{theorem}[Ferrero--Washington 1979, Ann.~Math. 109]
\label{v4-arith:pos:thm:FW}
Let \(K\) be an abelian number field, let \(K_\infty/K\)
be its cyclotomic \(\mathbb Z_p\)-extension, and let \(K_n\)
be its finite layers. With norm transition maps, put
\[
 X_K=\varprojlim_n\mathrm{Cl}(K_n)_p.
\]
Identify \(\operatorname{Gal}(K_\infty/K)\) with the chosen
\(\Gamma\). Then \(X_K\) is a finitely generated torsion
\(\Lambda\)-module and \(\mu(X_K)=0\).
\end{theorem}

\begin{remark}[applicability to Heisenberg bar]
\label{v4-arith:pos:rem:FW-applicability}
An application to an arithmetic Heisenberg bar must first specify
its finitely generated torsion coefficient module \(M\) and
its continuous \(\Gamma\)-action. It must then identify the
relevant localization with that of a module in
Theorem~\ref{v4-arith:pos:thm:FW}. Commuting creation operators
and a trivial action on an oscillator lattice supply no such
identification. The full Heisenberg Lie algebra itself has a
nonzero central commutator.

Section~\ref{v4-arith:pos:att-3-mu-heisenberg} proves vanishing for
finite free \(\mathbb Z_p\)-modules with trivial \(\Gamma\)-action
directly. The next statement gives a different, conditional
arithmetic route. Neither route supplies a conilpotence criterion.
\end{remark}

\begin{corollary}[transport of arithmetic \(\mu\)-vanishing]
\label{v4-arith:pos:cor:mu-zero-heisenberg}
Let \(M\) be a finitely generated torsion \(\Lambda\)-module
selected from a specified arithmetic bar construction.
Suppose \(M\) admits a pseudo-isomorphism to \(X_K^{\oplus d}\),
where \(K\) is abelian and \(X_K\) is as in
Theorem~\ref{v4-arith:pos:thm:FW}. Then \(\mu(M)=0\).
\end{corollary}

\begin{proof}
The finite kernel and cokernel have support at \((p,T)\).
They vanish after localization at \((p)\), giving
\[
 M_{(p)}\cong(X_K)_{(p)}^{\oplus d}=0.
\]
The last equality is Theorem~\ref{v4-arith:pos:thm:FW}.
Taking length proves the assertion. The pseudo-isomorphism is
an input, not a consequence of the name of the bar construction.
\end{proof}

\section{Constructing and Detecting the Comparison}
\label{v4-arith:pos:main-theorem}

\subsection{The Functors and Their Comparison Maps}
\label{v4-arith:pos:proof-overview}

The algebraic functors can be written before choosing a completed
recipient. Their descent to a derived category, compatibility with
reduction, and equivalence are separate assertions.
We first give the functors, then prove the completion criterion.
The algebraic bar--cobar adjunction and compatible quotient towers
provide two further comparisons.

\subsection{Step 1: The Algebraic Functors}
\label{v4-arith:pos:step-1}

Use the algebraic tensor product over \(\Lambda\) in this
subsection. Let \(C\) be a dg coalgebra and use its left and
right regular coactions. The contramodule and contratensor
definitions above extend by replacing \(k\) with \(\Lambda\).
A continuous version requires its own Hom and tensor objects.

\begin{definition}[the colinear Hom functor]
\label{v4-arith:pos:def:psi}
For a dg left comodule \(M\), put
\[
 \Psi_C(M)=\operatorname{Hom}_C(C,M)
 \subset\operatorname{Hom}_\Lambda(C,M),
 \qquad
 \operatorname{Hom}_\Lambda(C,M)^n
 =\prod_i\operatorname{Hom}_\Lambda(C^i,M^{i+n}).
\]
The differential is
\(d(f)=d_Mf-(-1)^{|f|}fd_C\).
The right coaction on \(C\) induces the left contraaction on
this colinear Hom complex. On morphisms, \(\Psi_C\) acts by
postcomposition.
\end{definition}

\begin{definition}[the contratensor functor]
\label{v4-arith:pos:def:phi}
For a dg left contramodule \(P\), put
\[
 \Phi_C(P)=C\odot_C P
 =\operatorname{coker}\bigl(
 C\otimes_\Lambda\operatorname{Hom}_\Lambda(C,P)
 \rightrightarrows C\otimes_\Lambda P\bigr).
\]
The two arrows are the contraaction and the right coaction
followed by signed evaluation. The differential and the left
coaction descend from \(C\otimes_\Lambda P\).
\end{definition}

Tensor--Hom adjunction identifies maps from this cokernel to
\(M\) with contramodule maps from \(P\) to
\(\operatorname{Hom}_C(C,M)\). Thus \(\Phi_C\dashv\Psi_C\)
on the algebraic dg categories. Completion is not part of
this calculation.

\begin{lemma}[descent through chosen resolution models]
\label{v4-arith:pos:lem:well-defined}
Suppose full triangulated resolution subcategories
\(\mathcal J\subset\mathrm{Hot}(C\text{-comod}_\Lambda)\)
and \(\mathcal P\subset\mathrm{Hot}(C\text{-contra}_\Lambda)\)
map equivalently to the chosen coderived and contraderived
recipients. If \(\Psi_C(\mathcal J)\subset\mathcal P\) and
\(\Phi_C(\mathcal P)\subset\mathcal J\), the restricted
adjunction transports to those recipients.
\end{lemma}
\begin{proof}
Compose each restricted functor with the two specified
resolution equivalences and their inverses.
The Hom adjunction transports along the same equivalences,
as do its unit, counit and triangle identities.
\end{proof}

Over a field, Theorem~\ref{v4-arith:pos:thm:classical-positselski}
supplies these resolution models and their equivalence.
Over \(\Lambda\), the models and their preservation must be
constructed. Injectivity over \(\Lambda\) alone does not prove
the lemma's hypotheses.

The contratensor quotient matters already in dimension two.
Over a field, let \(C=k\mathbf1\oplus kc\), with
\(\Delta c=\mathbf1\otimes c+c\otimes\mathbf1\).
The colinear endomorphisms of \(C\) have basis
\[
 f_0=\mathrm{id}_C,
 \qquad f_1(\mathbf1)=0,\quad f_1(c)=\mathbf1.
\]
Evaluation \(C\otimes\operatorname{Hom}_C(C,C)\to C\)
is surjective and has kernel
\[
 k(\mathbf1\otimes f_1)
 \oplus k(c\otimes f_1-\mathbf1\otimes f_0).
\]
These are precisely the contratensor relations, so the
quotient is \(C\). Omitting the quotient leaves a
four-dimensional space in place of a two-dimensional one.

\subsection{Step 2: Coefficient Change}
\label{v4-arith:pos:step-2}

\begin{lemma}[two coefficient comparisons]
\label{v4-arith:pos:lem:base-change}
Let \(R\to S\) be a map of commutative rings.
\begin{enumerate}
\item If \(C\) is a coaugmented conilpotent dg \(R\)-coalgebra,
then the algebraic scalar extension \(C\otimes_R S\)
is a coaugmented conilpotent dg \(S\)-coalgebra.
\item If \(E\) is a perfect complex of \(R\)-modules, then
for every \(M\in\mathrm D(R)\), the natural map
\[
 \operatorname{RHom}_R(E,M)\otimes_R^{\mathbb L}S
 \longrightarrow
 \operatorname{RHom}_S
   (E\otimes_R^{\mathbb L}S,M\otimes_R^{\mathbb L}S)
\]
is an equivalence.
\end{enumerate}
\end{lemma}
\begin{proof}
The coaugmentation splitting commutes with algebraic scalar
extension. Each element of its reduced part is a finite sum
of tensors \(c_i\otimes s_i\). A common bound on the vanishing
of the iterated reduced coproducts of the finitely many \(c_i\)
proves conilpotence. No flatness of \(S\) is needed for this
algebraic assertion.

For the second assertion, represent \(E\) by a bounded complex
of finite projective modules. Its Hom complex is
\(E^\vee\otimes_R M\), and
\(E^\vee\otimes_R S\cong(E\otimes_R S)^\vee\).
These identities give the stated derived map and its inverse.
\end{proof}

The second assertion concerns coefficient Hom.
The functor \(\Psi_C\) also imposes colinearity, and a
continuous version may involve an inverse limit.
Indeed, \(\operatorname{Hom}_\Lambda(\varinjlim C_i,M)
=\varprojlim\operatorname{Hom}_\Lambda(C_i,M)\).
Commutation of this limit, its derived functors and coefficient
change requires an additional proof.

For \(B=\Lambda/p\), a proposed comparison has the form
\begin{equation}
\label{v4-arith:pos:eq:base-change-diagram}
\begin{tikzcd}
 \mathcal C\arrow[r,"\Psi"]\arrow[d,"R_{\mathcal C}"']
 &\mathcal D\arrow[d,"R_{\mathcal D}"]\\
 \overline{\mathcal C}\arrow[r,"\overline\Psi"']
 &\overline{\mathcal D}.
\end{tikzcd}
\end{equation}
Here the barred recipients and the reduction functors must be
defined. A natural equivalence making this square commute is
part of the comparison. The corresponding square for \(\Phi\)
must also identify the units and counits.
An equivalence over \(B\), when supplied, is not the
field-coefficient theorem: \(B\) is not a field.

\subsection{Step 3: Detection by Derived Completion}
\label{v4-arith:pos:step-3}

\begin{proposition}[complete objects and reduction]
\label{v4-arith:pos:prop:lifting}
In a stable \(\Lambda\)-linear \(\infty\)-category with sequential limits,
a derived \(p\)-complete object \(X\) is zero if \(X/p=0\).
Consequently, a map whose cofiber is derived \(p\)-complete
is an equivalence if its reduction modulo \(p\) is.
\end{proposition}
\begin{proof}
The factorization \(p^n=p\,p^{n-1}\) gives a cofiber sequence
\[
 X/p^{n-1}\longrightarrow X/p^n\longrightarrow X/p.
\]
Induction gives \(X/p^n=0\) for every \(n\).
Derived completeness then gives
\(X\simeq\varprojlim_n X/p^n=0\).
Apply this result to the cofiber of the given map.
\end{proof}

This proves conservativity on the stated complete objects.
Faithfulness on morphisms fails.
Take \(M=\Lambda/(T)\cong\mathbb Z_p\).
It is derived \(p\)-complete and has \(\mu(M)=0\), yet
\[
 p\,\mathrm{id}_M\ne0,
 \qquad
 (p\,\mathrm{id}_M)\otimes_\Lambda^{\mathbb L}B=0.
\]
The two-term free resolution with differential \(T\)
verifies the second equality directly after reduction.

Lifts are not unique either.
The modules \(\Lambda/(T)\) and \(\Lambda/(T-p)\)
are \(\mathbb Z_p\)-free, derived \(p\)-complete, and have
\(\mu=0\). Their reductions are both \(B/(T)\).
They are not isomorphic over \(\Lambda\): \(T\) acts as zero
on the first and as multiplication by \(p\) on the second.
Even the first lift has the nonidentity automorphism
\((1+p)\mathrm{id}\) reducing to the identity.

\begin{corollary}[lifting a specified adjunction]
\label{v4-arith:pos:cor:psi-equivalence}
Consider an exact adjunction as in
Theorem~\ref{v4-arith:thm:arithmetic-positselski}.
Assume its unit and counit cofibers are derived \(p\)-complete
and its reduced unit and counit are equivalences.
Then the adjunction is an equivalence.
\end{corollary}
\begin{proof}
Proposition~\ref{v4-arith:pos:prop:lifting} makes the unit
and counit equivalences. The triangle identities identify
the two adjoint functors as inverses.
\end{proof}

An arithmetic application must verify these conditions in
its specified recipients. Applying ordinary reduction to
underlying complexes does not suffice unless that comparison
is itself conservative on the relevant cofibers.

\subsection{Step 4: The Algebraic Bar--Cobar Adjunction}
\label{v4-arith:pos:step-4}

Let \(A\) be an augmented dg \(\Lambda\)-algebra and \(C\)
a coaugmented conilpotent dg \(\Lambda\)-coalgebra.
Write \(\overline A=\ker\epsilon_A\) and
\(\overline C=\ker\epsilon_C\), and use cohomological shifts
\(V[1]^i=V^{i+1}\).
A twisting cochain is a degree-one map
\(\tau:\overline C\to\overline A\) satisfying
\[
 d_A\tau+\tau d_C+\mu_A(\tau\otimes\tau)\overline\Delta_C=0.
\]
The tensor of homogeneous maps includes its Koszul sign.

\begin{theorem}[algebraic \(\Lambda\)-bar--cobar adjunction]
\label{v4-arith:pos:thm:lambda-bar-cobar}
With algebraic tensors and direct sums, the constructions
\[
 B^\Lambda A=\bigoplus_{n\ge0}(\overline A[1])^{\otimes_\Lambda n},
 \qquad
 \Omega^\Lambda C=\bigoplus_{n\ge0}(\overline C[-1])^{\otimes_\Lambda n}
\]
carry their bar and cobar differentials and satisfy
\[
 \operatorname{Hom}_{\mathrm{dgAlg}^{\mathrm{aug}}_\Lambda}
       (\Omega^\Lambda C,A)
 \cong\operatorname{Tw}_\Lambda(C,A)
 \cong
 \operatorname{Hom}_{\mathrm{dgCoalg}^{\mathrm{conil}}_\Lambda}
       (C,B^\Lambda A).
\]
Thus \(\Omega^\Lambda\dashv B^\Lambda\) on these algebraic
dg categories.
\end{theorem}
\begin{proof}
An algebra map from the tensor algebra \(\Omega^\Lambda C\)
is determined by its restriction to \(\overline C[-1]\).
It commutes with differentials exactly when the corresponding
\(\tau\) satisfies the displayed twisting equation.

A coalgebra map to the tensor coalgebra \(B^\Lambda A\)
is determined by projection to \(\overline A[1]\).
Its length-\(n\) component applies that projection to the
\(n\)-fold reduced coproduct.
Conilpotence makes the resulting sum finite on each element.
The chain-map condition is the same twisting equation.
These two constructions are inverse and natural.
\end{proof}

No completion or derived equivalence is asserted by this
algebraic adjunction. A completed version must justify the
convergence of its sums and the weak equivalences it preserves.
It is also distinct from the module adjunction
\(\Phi_C\dashv\Psi_C\).

\subsection{Step 5: Unit and Counit}
\label{v4-arith:pos:step-5}

\begin{theorem}[a complete criterion for unit and counit]
\label{v4-arith:pos:thm:unit-counit-qi}
For the algebraic adjunction, the natural maps have types
\[
 \eta_C:C\longrightarrow B^\Lambda\Omega^\Lambda C,
 \qquad
 \varepsilon_A:\Omega^\Lambda B^\Lambda A\longrightarrow A.
\]
Suppose the cone of either map, viewed in \(\mathrm D(\Lambda)\),
is derived \(p\)-complete, and that its derived reduction
modulo \(p\) is a quasi-isomorphism. Then that map is a
quasi-isomorphism over \(\Lambda\).
\end{theorem}
\begin{proof}
Apply Proposition~\ref{v4-arith:pos:prop:lifting} in
\(\mathrm D(\Lambda)\) to the specified cone.
The maps exist by the algebraic adjunction. Their completeness
and reduction hypotheses must be checked for the chosen
\(A\) and \(C\).
\end{proof}

The coalgebra unit uses iterated \emph{reduced} coproducts.
The full coproduct of a primitive element retains terms with
\(\mathbf1\) in every arity, so conilpotence cannot truncate
that full coproduct. Nor does an equivalence in a chosen
coderived localization assert an ordinary quasi-isomorphism
without a comparison of the two notions.

\subsection{Step 6: Compatible Quotient Towers}
\label{v4-arith:pos:step-6}

\begin{lemma}[limits of compatible equivalences]
\label{v4-arith:pos:lem:independence}
Let \((\mathcal C_n)\) and \((\mathcal D_n)\) be inverse
diagrams of stable categories. A compatible equivalence of
diagrams \(\Psi_n:\mathcal C_n\simeq\mathcal D_n\) induces
\[
 \varprojlim_n\mathcal C_n\simeq\varprojlim_n\mathcal D_n.
\]
Restriction along a coinitial indexing functor preserves these
limits. In particular, cofinal systems of open ideals give
the same limit for a fixed diagram on those ideals.
\end{lemma}
\begin{proof}
The limit functor preserves equivalences of diagrams.
Its universal property also gives invariance under coinitial
restriction. For two systems of open ideals \(I_n,J_m\),
the quotients by \(I_n\cap J_m\) give a common refinement
with maps to the original quotients.
\end{proof}

The lemma requires the diagram and its compatibility maps.
Separate equivalences at individual quotients do not supply
a compatible equivalence of diagrams.
Identifying an arithmetic recipient with this inverse limit
is another hypothesis. The lemma gives no uniqueness for
arbitrary lifts of a residual object.

\subsection{The Comparison Criterion}
\label{v4-arith:pos:assembly}

\begin{proof}[Proof of Theorem~\ref{v4-arith:thm:arithmetic-positselski}]
Let \(K_P=\operatorname{cofib}(\eta_P)\) in \(\mathcal D\)
and \(L_M=\operatorname{cofib}(\varepsilon_M)\) in
\(\mathcal C\). They are derived \(p\)-complete by hypothesis.
Their reductions vanish because the reduced unit and counit
are equivalences. Proposition~\ref{v4-arith:pos:prop:lifting}
gives \(K_P=L_M=0\).
Thus both natural transformations are equivalences, and the
adjunction exhibits \(\Phi\) and \(\Psi\) as inverse functors.
\end{proof}

\begin{remark}[the dual arithmetic module]
\label{v4-arith:pos:rem:dual-mu}
A Selmer \(\mu\)-invariant requires a specified finitely
generated torsion Iwasawa module.
Comparing it with the invariant on another coefficient object
requires an arithmetic duality map and control of its kernel
and cokernel. A pseudo-isomorphism of the relevant modules
would identify their height-one localizations and hence their
\(\mu\)-invariants. Formal vector-space duality alone gives
no such arithmetic pseudo-isomorphism.
\end{remark}

\section{Application: Pairings and Arithmetic Levels}
\label{v4-arith:pos:heisenberg-application}

\subsection{The Required Arithmetic Comparison}
\label{v4-arith:pos:step-3-residual}

For Theorem~\ref{v4-arith:thm:P1-resolution} to apply to arithmetic
duality, the following assignments must be constructed:
\begin{itemize}
\item the arithmetic level \(\kappa^{\mathrm{Ar}}_{\mathsf G}\)
as a central invariant of a specified bar object.
\item the dual level \((\kappa^!)^{\mathrm{Ar}}_{\mathsf G}\)
as a central invariant of its specified dual object.
\item an involution and Mellin evaluation carrying their common
central defect to the difference of the two completed characters.
\end{itemize}
The converse also requires that evaluation to detect the defect.
A cobar algebra and a continuous dual algebra have different
constructions. A comparison with a Bost--Connes algebra must specify
which is used. The local weight calculation below proves evaluation
compatibility for its oscillator model. It supplies none of these
arithmetic identifications.

\subsection{The Explicit Weight Computation}
\label{v4-arith:pos:weight-comp}

\begin{theorem}[evaluation pairing and arithmetic energy]
\label{v4-arith:pos:thm:weight-comp}
For a finite set \(S\) of arithmetic primes let
\(P_S=\mathbb Z[x_{p,k}:p\in S,k\ge1]\), with oscillator weight
\(k\) and energy \(k\log p\) on \(x_{p,k}\).
Each weight component \(P_{S,w}\) is finite free. Its restricted dual
\[
 P_S^\circ=\bigoplus_{w\ge0}\operatorname{Hom}_{\mathbb Z}
                              (P_{S,w},\mathbb Z)
\]
has a perfect evaluation pairing in each weight. Give the dual of a
monomial energy equal to the negative of that monomial's energy.
Then evaluation has total energy zero.

If an arithmetic bar--cobar pairing and comparison maps identify these
particular weight components and their pairings with \(P_{S,w}\) and
its restricted dual, the same energy cancellation holds on their images.
An identification of a dual generator with \(\mu_{p^k}\), and an
identification of this energy cancellation with
\(\kappa^{\mathrm{Ar}}_{\mathsf G}+(\kappa^!)^{\mathrm{Ar}}_{\mathsf G}=0\),
are separate comparison hypotheses.
\end{theorem}
\begin{proof}
There are finitely many monomials of a given oscillator weight because
\(S\) is finite. Write \(x^\nu\) for their monomial basis and
\(\eta_\nu\) for the dual basis. Evaluation is
\(\eta_\mu(x^\nu)=\delta_{\mu\nu}\), hence perfect in every weight.
Its only nonzero basis pairings have energies
\(E(\nu)\) and \(-E(\nu)\), whose sum is zero. Transport through the
specified pairing-preserving maps proves the conditional assertion.
No arithmetic map is constructed by this transport calculation.

The distinction from a polynomial generator algebra can already be seen
in one variable: dualizing the primitive coproduct of \(\mathbb Z[x]\)
gives \(\eta_a\eta_b=\binom{a+b}{a}\eta_{a+b}\), so
\(\eta_1^k=k!\eta_k\). This follows by evaluating on
\(\Delta x^{a+b}=\sum_j\binom{a+b}{j}x^j\otimes x^{a+b-j}\).
It is not the relation \(\mu_{p^k}=\mu_p^k\) with the same integral
basis. A different normalization or cobar differential requires its own
construction.
\end{proof}

\subsection{Global Euler-Product Check}
\label{v4-arith:pos:euler-check}

\begin{proposition}[global weight sum]
\label{v4-arith:pos:prop:global-weight}
The formal cancellation of opposite energy labels over prime-power
indices is
\begin{equation}
\sum_{p\text{ prime}}\sum_{n\geq 1}\bigl[n\log p\cdot p^{-ns}+(-n\log p)\cdot p^{-ns}\bigr]\;=\;0,
\end{equation}
convergent absolutely for \(\mathrm{Re}(s)>1\). After termwise
pairing it is the zero function, hence extends as zero. No
meromorphic continuation of the unpaired prime-indexed sum is asserted.
\end{proposition}

\begin{proof}
Term-by-term: at each prime \(p\) and each \(n\geq 1\), the two
weights \(n\log p\) and \(-n\log p\) sum to zero. For \(\sigma=\operatorname{Re}s>1\), the unpaired absolute sum is
\[
 \sum_p(\log p)\frac{p^{-\sigma}}{(1-p^{-\sigma})^2}
 \le (1-2^{-\sigma})^{-2}\sum_{m\ge2}(\log m)m^{-\sigma}<\infty.
\]
Thus pairing terms is justified in that half-plane, and their sum is
zero. This scalar identity supplies no arithmetic pairing or level map. The identity
\(-\zeta'(s)/\zeta(s)=\sum_{p,k\geq1}(\log p)p^{-ks}\) is a prime-power
identity, not an identity for the prime-indexed sum alone.
\end{proof}

\begin{remark}[pairing and level comparison]
\label{v4-arith:pos:cor:step-3-unconditional}
The local evaluation theorem proves energy cancellation. To use it in
Theorem~\ref{v4-arith:thm:P1-resolution}, one must construct the
arithmetic pairing, the central-defect identity, and the Mellin
evaluation, including its detection of the defect. Opposite energy
labels alone do not identify these central classes.
The five conditions of Remark~\ref{v4-arith:pos:rem:open-sub-items} must be
realized with those compatibility maps; a named categorical equivalence
alone does not identify its numerical invariants.
\end{remark}

\subsection{The Arithmetic Comparison Required for P1}
\label{v4-arith:pos:p1-closes}
\begin{question}[level cancellation and the functional equation]
\label{v4-arith:pos:thm:P1-full-unconditional}
Construct the arithmetic pairing, its integral and completed coefficient
models, and the comparison of its central obstruction classes with the
Tate functional equation. The intended conclusion remains
\[
 \kappa^{\mathrm{Ar}}_{\mathsf G}+(\kappa^!)^{\mathrm{Ar}}_{\mathsf G}=0
 \quad\Longleftrightarrow\quad \xi(s)=\xi(1-s).
\]
The scalar functional equation is a theorem. The oscillator energy
calculation does not prove either identification of the arithmetic
classes with its two sides.
\end{question}

\begin{remark}[dependencies of the arithmetic comparison]
\label{v4-arith:pos:rem:form-prop}
The pairing and invariant comparisons are required whenever
Remark~\ref{v4-arith:cl:rem:P1-conditional} and Question~\ref{v4-arith:thm:G-row-complementarity} and Theorem~\ref{v4-arith:rh:thm:sufficient-path} invoke the arithmetic Positselski
input. The continuous-dual, completion, lifting, character and
Bost--Connes requirements remain distinct. Neither the formal sum of
opposite energies nor the scalar zeta identity removes these maps.
\end{remark}

\section{Fourier Pairings and Adelic Assembly}
\label{v4-arith:pos:archimedean}

\subsection{The Archimedean Fourier Model}
\label{v4-arith:pos:archimedean-analogue}

The real local factor enters the scalar functional equation
through the Gaussian. Fix the Fourier convention
\[
 \widehat f(y)=\int_{\mathbb R}f(x)e^{-2\pi ixy}\,dx.
\]
This integral uses real analysis and a specified measure.
A comparison with an algebraic pairing must preserve that
normalization.

\begin{theorem}[Gaussian Fourier identity]
\label{v4-arith:pos:thm:gaussian-positselski}
For \(t>0\) and \(f_t(x)=e^{-\pi t x^2}\),
\[
 \widehat f_t(y)=t^{-1/2}e^{-\pi y^2/t}.
\]
\end{theorem}
\begin{proof}
For \(t=1\), differentiate under the integral and integrate
by parts. The Gaussian decay eliminates the boundary term
and gives \(\widehat f_1'(y)=-2\pi y\widehat f_1(y)\).
The value at zero is one: the square of the positive Gaussian
integral is \(2\pi\int_0^\infty r e^{-\pi r^2}\,dr=1\).
Solving the differential equation gives the formula at \(t=1\).
The substitution \(u=\sqrt t\,x\) proves the general case.
\end{proof}

Poisson summation gives
\(\theta(t)=t^{-1/2}\theta(1/t)\) for
\(\theta(t)=\sum_{m\in\mathbb Z}e^{-\pi m^2t}\).
The theta series counts a lattice quadratic form.
The oscillator Fock space
\(\mathbb C[\alpha_{-1},\alpha_{-2},\ldots]|0\rangle\),
with weight \(n\) on \(\alpha_{-n}\), has character
\(\prod_{n\ge1}(1-q^n)^{-1}\).
The two series already differ at the coefficient of \(q\):
\(\sum_{m\in\mathbb Z}q^{m^2}\) has coefficient two,
whereas the oscillator character has coefficient one.

To compare the Fourier identity with co--contra duality,
specify an archimedean coalgebra and a pairing-preserving map.
Theorem~\ref{v4-arith:pos:thm:classical-positselski}
applies to an actual dg coalgebra over \(\mathbb C\).
It does not turn a Fock module, a theta series or a Fourier
transform into that coalgebra or its categorical pairing.

\subsection{Adelic Assembly}
\label{v4-arith:pos:adelic-assembly}

Fix a commutative coefficient ring \(R\).
For each place \(v\), let \(\mathcal C_v,\mathcal D_v\)
be presentable stable \(R\)-linear categories, with chosen
objects \(c_v,d_v\). Use the tensor product of such categories
characterized by \(R\)-bilinear functors preserving colimits
in each variable.

For a finite place set \(F\), put
\[
 \mathcal C_F=\bigotimes_{v\in F,R}\mathcal C_v,
 \qquad
 \mathcal D_F=\bigotimes_{v\in F,R}\mathcal D_v.
\]
The empty tensor is \(\mathrm{Mod}_R\).
Insertion of \(c_v\), respectively \(d_v\), defines the
transition functors when a new place is added.
Define the corresponding restricted tensor categories as
colimits over these finite place sets.

\begin{theorem}[assembly of pointed local equivalences]
\label{v4-arith:pos:thm:adelic-positselski}
Suppose \(R\)-linear equivalences
\(\Psi_v:\mathcal C_v\simeq\mathcal D_v\) and identifications
\(\Psi_v(c_v)\simeq d_v\) give a compatible equivalence of
the finite tensor diagrams. Then
\begin{equation}
\label{v4-arith:pos:eq:adelic-positselski}
 \varinjlim_F\mathcal C_F
 \xrightarrow{\ \sim\ }
 \varinjlim_F\mathcal D_F.
\end{equation}
\end{theorem}
\begin{proof}
Tensoring finitely many equivalences gives an equivalence
\(\mathcal C_F\simeq\mathcal D_F\).
The pointing identifications and their assumed coherence make
these maps an equivalence of diagrams. Taking the colimit
preserves that equivalence.
\end{proof}

For arithmetic local categories, the common coefficient action,
tensor recipient, local equivalences and compatible pointings
must be constructed. The symbols \(\mathbb Q_p\) and
\(\mathbb C\) do not provide a common field by themselves.
The chosen objects also need not be scalar augmentations or
vacua of an independently specified field theory.
Identifying this categorical colimit with an analytic adelic
restricted tensor product is a further comparison.

\begin{remark}[the adelic pairing and the functional equation]
\label{v4-arith:pos:rem:adelic-fe}
The theta--Mellin calculation of
Theorem~\ref{v4-arith:thm:arithmetic-modular-invariance} retains
its polar terms and proves the scalar functional equation.
To interpret it through an adelic bar pairing, construct an
evaluation identifying the duality involution with Mellin reflection.
The forward implication also requires the central-defect identity.
The converse requires that evaluation to detect the specified defect,
as in Theorem~\ref{v4-arith:thm:P1-resolution}.

An adelic co--contra equivalence is a separate categorical
comparison. Identifying its pairing with Fourier self-duality
requires the actual coefficient and pairing maps.
The scalar theta calculation does not construct them.
\end{remark}

\section{Residual Conditions}
\label{v4-arith:pos:residual-domain}

The integral lifting criterion and the local assembly theorem
apply only after their functors, recipients and comparison
maps have been supplied. The following questions ask for further arithmetic structures.
The commutative comparison remains subject to the preceding
hypotheses.

\subsection{Non-Abelian Iwahori Level}
\label{v4-arith:pos:non-abelian}

For a compact \(p\)-adic group \(G\), write
\(A=\mathbb Z_p[[G]]\). An \(A\)-coring is an
\(A\)-bimodule \(C\) with counit \(C\to A\) and a
coassociative coproduct \(C\to C\otimes_A C\).
A continuous version replaces the tensor by a specified
completed tensor and requires continuous structure maps.

\begin{question}[noncommutative arithmetic co--contra comparison]
\label{v4-arith:pos:cj:non-commutative}
For an arithmetic coring over \(A\), construct exact comodule
and contramodule recipients and an adjunction between them.
Determine a coefficient and completion domain on which its
unit and counit are equivalences.
For \(G=GL_2(\mathbb Z_p)\), construct the required comparison
with the intended Iwahori Hecke category and its actions.
If a noncommutative \(\mu\)-invariant is used, specify its
module and prove the precise role of its vanishing.
\end{question}

\begin{remark}[module structure and categorical comparison]
\label{v4-arith:pos:rem:non-commutative-form}
A noncommutative module structure theorem has its own hypotheses
on the group, ring and torsion class.
It does not by itself compare comodule and contramodule
resolutions. Such a comparison needs the actual tensor and
Hom functors, with the correct module sides.
Even over the commutative ring \(\Lambda\),
Proposition~\ref{v4-arith:pos:prop:resolutions-cohere}
shows why \(\mu=0\) cannot replace the necessary Tor conditions.
\end{remark}

\subsection{Whittaker Model Compatibility}
\label{v4-arith:pos:whittaker}

A Whittaker condition requires a specified unipotent group
\(N\), a character \(\psi\), and an action on the recipient.
It also requires a choice of equivariant or derived
coinvariant construction. The coefficient ring must support
that character.

\begin{question}[Whittaker-compatible co--contra comparison]
\label{v4-arith:pos:cj:whittaker-compat}
Construct an arithmetic co--contra comparison equipped with
compatible \((N,\psi)\)-actions and the required Hecke actions.
Show that it induces an equivalence between the resulting
Whittaker recipients,
\[
 \mathcal C^{N,\psi}\simeq\mathcal D^{N,\psi},
\]
and construct its comparison with the chosen Satake and
Whittaker functors. A scalar character identity alone does
not define these equivariant functors.
\end{question}

\subsection{Extension to Higher Rank}
\label{v4-arith:pos:higher-rank}

Replacing \(\Gamma\cong\mathbb Z_p\) by \(\mathbb Z_p^d\)
changes the coefficient ring to
\(\Lambda_d=\mathbb Z_p[[T_1,\ldots,T_d]]\).
The co--contra comparison then needs categories and resolution
hypotheses over this ring. An arithmetic vanishing theorem also
requires a specified module attached to the chosen tower.
The cyclotomic theorem does not supply that module or an induction
in \(d\).

\begin{remark}[higher-rank extension]
\label{v4-arith:pos:rem:higher-rank}
For a finitely generated torsion \(\Lambda_d\)-module, the
length of its localization at \((p)\) gives a corresponding
height-one invariant. The module \(\Lambda_d/(p)\) has length
one there, although \(\Gamma\) is abelian. Thus abelianity alone
does not force this invariant to vanish. Determining its value
for an arithmetic tower and comparing that module with a bar
coefficient are separate problems.
\end{remark}

\section{Positselski Boundary}
\label{v4-arith:pos:summary}

\begin{enumerate}
\item Theorem~\ref{v4-arith:thm:arithmetic-positselski} gives
a lifting criterion for a specified adjunction.
Its unit and counit cofibers must be derived \(p\)-complete,
and their reductions must vanish.
The integral arithmetic categories and functors remain
additional constructions. The field correspondence does
not establish them by scalar extension or faithful reduction.
\item The comparisons of Section~\ref{v4-arith:pos:obstruction-analysis-table}
have separate hypotheses. In particular, the weight calculation needs
arithmetic pairing and invariant maps even in the abelian case.
The nonabelian Iwahori extension is a further problem
(Section~\ref{v4-arith:pos:residual-domain}).
\item The hypothesis \(\mu(C)=0\) concerns the specified torsion
coefficient module. Section~\ref{v4-arith:pos:att-3-mu-heisenberg}
proves a local lattice calculation.
Corollary~\ref{v4-arith:pos:cor:mu-zero-heisenberg} gives an
arithmetic consequence under an actual module comparison.
Neither statement identifies an unspecified Heisenberg bar with
a cyclotomic class-group module or characterizes conilpotence.
\item Theorem~\ref{v4-arith:pos:thm:weight-comp} constructs the restricted
integral dual and its energy-zero evaluation pairing on finite-species
polynomial components. Identifying this pairing with arithmetic
bar--cobar data and its energies with central obstruction classes is
additional work.
\item The intended P1 comparison remains
Question~\ref{v4-arith:pos:thm:P1-full-unconditional}. Its operator and invariant
maps are not consequences of the scalar functional equation.
\item Adelic extension retains the local categorical hypotheses
and their compatible pointings. Its interpretation through the
functional equation additionally requires the evaluation and
central-defect maps of Remark~\ref{v4-arith:pos:rem:adelic-fe}.
Gaussian Fourier self-duality and the scalar theta calculation
provide analytic models, not those arithmetic identifications.
\item Residual condition
(Section~\ref{v4-arith:pos:residual-domain}): non-commutative arithmetic
Positselski for non-abelian Iwahori level
(Question~\ref{v4-arith:pos:cj:non-commutative}), Whittaker compatibility
(Question~\ref{v4-arith:pos:cj:whittaker-compat}) and higher-rank
Iwasawa coefficients (Remark~\ref{v4-arith:pos:rem:higher-rank}).
\end{enumerate}

The intended arithmetic comparison relates central obstruction classes
to the classical functional equation. The local constructions above do
not identify those classes with the explicitly computed scalar pairing.
The continuous-dual, completion, lifting, character and Bost--Connes
requirements must therefore be supplied as actual compatible maps on
the intended carrier. The arithmetic Positselski input on the sufficient
path of Theorem~\ref{v4-arith:rh:thm:sufficient-path} includes that comparison;
it is not discharged by the equality of opposite energy labels.
